\chapter{Desarrollo de Aplicaciones Móviles Inclusivas}
\label{cp:aplicaciones-moviles}

\insertminitoc
\parindent0pt

Una aplicación de reconocimiento de lengua de señas para niños sordos no es únicamente
un proyecto de inteligencia artificial; es, ante todo, un producto de software que debe
funcionar con fiabilidad en un dispositivo móvil, procesar video en tiempo real, presentar
una interfaz comprensible para un público infantil y enmarcarse dentro del campo de la
comunicación aumentativa y alternativa. El presente capítulo aborda estos cuatro ejes: el
framework de desarrollo multiplataforma utilizado, la integración de modelos de
inteligencia artificial con la cámara del dispositivo, los principios de diseño de
interfaces accesibles para niños y el marco conceptual de las tecnologías de comunicación
aumentativa y alternativa (\acrshort{caa}) en el que se inserta el sistema.

\section{Desarrollo Móvil Multiplataforma con Flutter}

Flutter es un framework de código abierto desarrollado por Google para la construcción de
aplicaciones compiladas de forma nativa a partir de una base de código única en el lenguaje
Dart \cite{google_flutter_2024}. A diferencia de otros frameworks multiplataforma que
dependen de un puente de comunicación con los componentes nativos del sistema operativo,
Flutter compila el código Dart de forma anticipada (ahead-of-time, \acrshort{aot}) a código de
máquina nativo y utiliza su propio motor de renderizado gráfico para dibujar cada píxel de
la interfaz, lo que produce un rendimiento cercano al de una aplicación nativa
\cite{google_flutter_2024}.

La arquitectura de Flutter se organiza en capas. La capa inferior es el motor (engine),
escrito en C++, que se encarga del renderizado gráfico, la composición, la gestión de
entrada y salida y la comunicación con el sistema operativo. Sobre esta capa se encuentra
el framework en Dart, que proporciona un sistema de widgets reactivos con el cual se
construye toda la interfaz de usuario. Cada elemento visual, desde un botón hasta la
distribución de la pantalla, es un widget que se combina de forma declarativa con otros
widgets para formar interfaces complejas \cite{google_flutter_2024}. Este modelo
declarativo facilita la construcción de interfaces dinámicas que se actualizan en respuesta
a cambios de estado, lo que resulta especialmente útil para una aplicación que muestra
resultados de clasificación en tiempo real. Flutter permite generar aplicaciones para
Android e iOS desde el mismo código, lo cual reduce el esfuerzo de desarrollo y
mantenimiento del sistema.

\section{Integración de Modelos de Inteligencia Artificial y Procesamiento de Cámara en Tiempo Real}

El sistema propuesto requiere que la cámara frontal del dispositivo capture cuadros de
video de forma continua, que MediaPipe Hands y MediaPipe Pose extraigan los landmarks de
las manos y del cuerpo en cada cuadro, y que los modelos descritos en el capítulo 6, la
red convolucional 1D para el alfabeto y la red \acrshort{lstm} para las señas dinámicas, clasifiquen
los gestos detectados. Todo este procesamiento debe ocurrir en tiempo real, sin bloquear
la interfaz de usuario y sin enviar datos a un servidor externo.

En Flutter, el acceso a la cámara se gestiona a través del plugin camera, que proporciona
un flujo de cuadros (stream) desde la cámara del dispositivo. Cada cuadro se transfiere
como un arreglo de bytes al canal de procesamiento de MediaPipe Hands, que devuelve las
coordenadas de los 21 landmarks por mano detectada \cite{google_mediapipe_2024},
\cite{zhang_mediapipe_2020}. Estos landmarks se normalizan y se pasan al intérprete de
TensorFlow Lite para obtener la predicción del modelo \cite{google_tflite_2024}. Para
evitar que este procesamiento bloquee el hilo principal de la interfaz, la extracción de
landmarks con MediaPipe se ejecuta en un hilo nativo en segundo plano, comunicado con
Flutter mediante un canal de plataforma, de modo que el hilo de la interfaz permanece libre. Sousa et al. validaron la viabilidad de esta arquitectura en un
sistema de reconocimiento de Lengua de Señas Portuguesa implementado con Flutter,
MediaPipe y \acrshort{tflite}, donde reportaron tiempos de inferencia compatibles con una
experiencia de uso fluida en dispositivos Android de gama media \cite{sousa_lgp_flutter_2025}.

\section{Diseño de Interfaces para Niños y Accesibilidad}

El público objetivo del sistema es la población infantil sorda, lo que impone restricciones
de diseño particulares. Hourcade señala que las interfaces dirigidas a niños deben
minimizar el uso de texto, privilegiar representaciones visuales e iconográficas, utilizar
elementos de interacción de tamaño suficiente para las habilidades motoras en desarrollo y
ofrecer retroalimentación inmediata ante cada acción del usuario \cite{hourcade_children_2008}.
En el caso de niños sordos, estas consideraciones se intensifican: la lengua escrita suele
ser una segunda lengua adquirida con dificultad, por lo que la interfaz debe apoyarse en
recursos visuales como colores, iconos y animaciones para comunicar el estado del sistema
y los resultados de la clasificación.

Las directrices de accesibilidad de contenido web \acrshort{wcag} 2.1, publicadas por el \acrshort{w3c} en
2018, establecen cuatro principios fundamentales: el contenido debe ser perceptible,
operable, comprensible y robusto \cite{w3c_wcag_2018}. Si bien estas directrices fueron
concebidas para contenido web, sus principios son aplicables al diseño de interfaces
móviles. En particular, el principio de perceptibilidad exige que la información se
presente de formas que el usuario pueda percibir, lo que para un usuario sordo implica que
toda la retroalimentación del sistema debe ser visual o háptica, nunca exclusivamente
auditiva. Flutter proporciona soporte nativo para herramientas de accesibilidad del sistema
operativo como TalkBack (Android) y VoiceOver (iOS), y permite configurar etiquetas
semánticas en los widgets para que los lectores de pantalla describan correctamente cada
elemento de la interfaz \cite{google_flutter_2024}. La Convención sobre los Derechos de
las Personas con Discapacidad de las Naciones Unidas refuerza este enfoque al establecer
que las tecnologías de la información deben ser accesibles y estar disponibles para las
personas con discapacidad en igualdad de condiciones \cite{onu_cdpd_2006}.

\section{Tecnologías de Comunicación Aumentativa y Alternativa}

El sistema desarrollado no es únicamente una aplicación móvil con modelos de inteligencia
artificial; desde el punto de vista comunicativo, se enmarca dentro de un campo más amplio
con una tradición consolidada de herramientas diseñadas para apoyar la comunicación de
personas que no pueden depender exclusivamente del habla oral.

La comunicación aumentativa y alternativa (\acrshort{caa}) es un término general que abarca todos
los métodos de comunicación distintos al habla oral que permiten a las personas expresar
pensamientos, necesidades y sentimientos. La American Speech-Language-Hearing Association
(\acrshort{asha}) define la \acrshort{caa} como un conjunto de herramientas y estrategias que incluye tableros
de imágenes, dispositivos generadores de habla, signos manuales, gestos y deletreo con los
dedos \cite{asha_aac_nodate}. Los sistemas de \acrshort{caa} se clasifican en dos categorías: los
sistemas sin ayuda (unaided), que utilizan únicamente el cuerpo del usuario, como las
lenguas de señas y los gestos, y los sistemas con ayuda (aided), que requieren un
dispositivo externo, desde un tablero de imágenes impreso hasta una aplicación en un
teléfono inteligente \cite{beukelman_aac_2020}.

Beukelman y Light clasifican los sistemas con ayuda según su nivel tecnológico: sin
tecnología (tableros de papel), baja tecnología (dispositivos con botones y mensajes
pregrabados) y alta tecnología (software con generación de habla, predicción de texto y
acceso a vocabulario extenso) \cite{beukelman_aac_2020}. La aplicación desarrollada en
el presente proyecto se ubica en esta última categoría, con la particularidad de que
integra visión por computadora para interpretar los gestos del usuario en lugar de
depender de la selección manual de símbolos en pantalla. La combinación de un módulo de
reconocimiento de señas (dirección niño a oyente) con un diccionario visual de señas
(dirección oyente a niño) convierte al sistema en una herramienta de
comunicación bidireccional que abarca tanto la función aumentativa como la alternativa de
la \acrshort{caa}, adaptada específicamente a las necesidades de la comunidad sorda hondureña y a la
Lengua de Señas Hondureña (\acrshort{lesho}).
